\documentclass[10pt,a4paper]{amsart}
\usepackage{amsmath,amssymb,amsthm,graphicx}
\usepackage{verbatim}
\newenvironment{reference}{\comment}{\endcomment}
%\newenvironment{reference}{}{}
\newenvironment{slogan}{\comment}{\endcomment}
\newenvironment{history}{\comment}{\endcomment}

% For commutative diagrams we use Xy-pic
\usepackage[all]{xy}

% We use 2cell for 2-commutative diagrams.
\xyoption{2cell}
% Theorem environments.
%
\theoremstyle{plain}
\newtheorem{theorem}[subsection]{Theorem}
\newtheorem{proposition}[subsection]{Proposition}
\newtheorem{lemma}[subsection]{Lemma}

\theoremstyle{definition}
\newtheorem{definition}[subsection]{Definition}
\newtheorem{example}[subsection]{Example}
\newtheorem{exercise}[subsection]{Exercise}
\newtheorem{situation}[subsection]{Situation}

\theoremstyle{remark}
\newtheorem{remark}[subsection]{Remark}
\newtheorem{remarks}[subsection]{Remarks}

\numberwithin{equation}{subsection}

% Macros
%
\def\lim{\mathop{\mathrm{lim}}\nolimits}
\def\colim{\mathop{\mathrm{colim}}\nolimits}
\def\Spec{\mathop{\mathrm{Spec}}}
\def\Hom{\mathop{\mathrm{Hom}}\nolimits}
\def\Ext{\mathop{\mathrm{Ext}}\nolimits}
\def\SheafHom{\mathop{\mathcal{H}\!\mathit{om}}\nolimits}
\def\SheafExt{\mathop{\mathcal{E}\!\mathit{xt}}\nolimits}
\def\Sch{\mathit{Sch}}
\def\Mor{\mathop{\mathrm{Mor}}\nolimits}
\def\Ob{\mathop{\mathrm{Ob}}\nolimits}
\def\Sh{\mathop{\mathit{Sh}}\nolimits}
\def\NL{\mathop{N\!L}\nolimits}
\def\CH{\mathop{\mathrm{CH}}\nolimits}
\def\proetale{{pro\text{-}\acute{e}tale}}
\def\etale{{\acute{e}tale}}
\def\QCoh{\mathit{QCoh}}
\def\Ker{\mathop{\mathrm{Ker}}}
\def\Im{\mathop{\mathrm{Im}}}
\def\Coker{\mathop{\mathrm{Coker}}}
\def\Coim{\mathop{\mathrm{Coim}}}

% Boxtimes
%
\DeclareMathSymbol{\boxtimes}{\mathbin}{AMSa}{"02}

%
% Macros for moduli stacks/spaces
%
\def\QCohstack{\mathcal{QC}\!\mathit{oh}}
\def\Cohstack{\mathcal{C}\!\mathit{oh}}
\def\Spacesstack{\mathcal{S}\!\mathit{paces}}
\def\Quotfunctor{\mathrm{Quot}}
\def\Hilbfunctor{\mathrm{Hilb}}
\def\Curvesstack{\mathcal{C}\!\mathit{urves}}
\def\Polarizedstack{\mathcal{P}\!\mathit{olarized}}
\def\Complexesstack{\mathcal{C}\!\mathit{omplexes}}
% \Pic is the operator that assigns to X its picard group, usage \Pic(X)
% \Picardstack_{X/B} denotes the Picard stack of X over B
% \Picardfunctor_{X/B} denotes the Picard functor of X over B
\def\Pic{\mathop{\mathrm{Pic}}\nolimits}
\def\Picardstack{\mathcal{P}\!\mathit{ic}}
\def\Picardfunctor{\mathrm{Pic}}
\def\Deformationcategory{\mathcal{D}\!\mathit{ef}}
\usepackage{cleveref}
\crefname{originallemma}{original lemma}{original lemmas}
\Crefname{originallemma}{Original lemma}{Original lemmas}
\makeatletter
\newlabel{lemma-base-change-quasi-finite}{{03XL}{source}{}{sourcepointer-1}{}}
\newlabel{lemma-base-change-quasi-finite@cref}{{[originallemma][0][]03XL}{[1][0][]source}{}{}{}}
\makeatother
\crefname{originallemma}{original lemma}{original lemmas}
\Crefname{originallemma}{Original lemma}{Original lemmas}
\makeatletter
\newlabel{morphisms-lemma-dimension-fibre-specialization}{{06RU}{source}{}{sourcepointer-2}{}}
\newlabel{morphisms-lemma-dimension-fibre-specialization@cref}{{[originallemma][0][]06RU}{[1][0][]source}{}{}{}}
\makeatother
\crefname{originallemma}{original lemma}{original lemmas}
\Crefname{originallemma}{Original lemma}{Original lemmas}
\makeatletter
\newlabel{morphisms-lemma-etale-over-field}{{02GL}{source}{}{sourcepointer-3}{}}
\newlabel{morphisms-lemma-etale-over-field@cref}{{[originallemma][0][]02GL}{[1][0][]source}{}{}{}}
\makeatother
\crefname{originallemma}{original lemma}{original lemmas}
\Crefname{originallemma}{Original lemma}{Original lemmas}
\makeatletter
\newlabel{morphisms-lemma-locally-quasi-finite-fibres}{{06RT}{source}{}{sourcepointer-4}{}}
\newlabel{morphisms-lemma-locally-quasi-finite-fibres@cref}{{[originallemma][0][]06RT}{[1][0][]source}{}{}{}}
\makeatother
\crefname{originallemma}{original lemma}{original lemmas}
\Crefname{originallemma}{Original lemma}{Original lemmas}
\makeatletter
\newlabel{morphisms-lemma-quasi-finite-at-point-characterize}{{01TH}{source}{}{sourcepointer-5}{}}
\newlabel{morphisms-lemma-quasi-finite-at-point-characterize@cref}{{[originallemma][0][]01TH}{[1][0][]source}{}{}{}}
\makeatother
\crefname{originallemma}{original lemma}{original lemmas}
\Crefname{originallemma}{Original lemma}{Original lemmas}
\makeatletter
\newlabel{spaces-properties-lemma-base-change-etale}{{03FU}{source}{}{sourcepointer-6}{}}
\newlabel{spaces-properties-lemma-base-change-etale@cref}{{[originallemma][0][]03FU}{[1][0][]source}{}{}{}}
\makeatother
\crefname{originallemma}{original lemma}{original lemmas}
\Crefname{originallemma}{Original lemma}{Original lemmas}
\makeatletter
\newlabel{spaces-properties-lemma-points-cartesian}{{03H4}{source}{}{sourcepointer-7}{}}
\newlabel{spaces-properties-lemma-points-cartesian@cref}{{[originallemma][0][]03H4}{[1][0][]source}{}{}{}}
\makeatother
\title{A locally finite type morphism of algebraic spaces is locally quasi-finite exactly when its field-valued fibres are discrete}
\author{The Stacks Project Authors}
\date{stacks-project:spaces-morphisms.tex:local-snapshot-20261012; GNU FDL 1.2 or later}
\begin{document}
\maketitle
\setcounter{subsection}{1}
\noindent\textit{Editorial note.} All delivered mathematical text is the author's own, copied from the pinned Stacks Project chapter file spaces-morphisms.tex; no statement or proof step was rewritten, reordered or removed.\par\medskip
\noindent\textit{Editorial note.} The proof is the author's own two-way argument; the converse direction builds the contradiction from the points of the etale equivalence relation R\_v rather than citing the scheme case.\par\medskip
\noindent\textit{Editorial note.} The seven references of the counted spans are recorded as non-load-bearing navigation pointers; the argument that consumes them stays inside the counted proof span.\par\medskip
\noindent\textit{Editorial note.} Preamble: the three native excerpt spans of sources/stacks-spaces-morphisms/preamble.tex (lines 8-13, 14-18 and 154-217). The chapter file itself carries no preamble text.\par\medskip
\noindent\textit{Allowed standard prerequisites (editorial).}\begin{itemize}
\item Base change of etale morphisms of algebraic spaces. Properties of Spaces, Lemma spaces-properties-lemma-base-change-etale, used at proof lines 5420-5422. Makes U\_k -> X\_k etale, hence open and surjective on points.
\item Stability of local quasi-finiteness under base change. Lemma lemma-base-change-quasi-finite, used at proof lines 5423-5425. Transfers local quasi-finiteness to the fibre X\_k.
\item The scheme criterion: local quasi-finiteness at a point via discrete fibres. Morphisms, Lemmas morphisms-lemma-locally-quasi-finite-fibres and morphisms-lemma-quasi-finite-at-point-characterize, used at proof lines 5427-5428 and 5444-5447. Gives discreteness in one direction and produces the specialization in the other.
\item Dimension formula for specialisations of points of a scheme over a field. Morphisms, Lemma morphisms-lemma-dimension-fibre-specialization, used at proof lines 5451-5454. Yields the strict inequality of transcendence degrees.
\item Residue fields of etale morphisms over a field are finite separable extensions. Morphisms, Lemma morphisms-lemma-etale-over-field, and Properties of Spaces, Lemma spaces-properties-lemma-points-cartesian, used at proof lines 5458-5464. Forces the two transcendence degrees to be equal.
\end{itemize}\medskip
\ifdefined\hypertarget\hypertarget{sourcepointer-1}{}\fi
\noindent\textit{Source reference.} Original lemma 03XL: focused-seeds/AGP/sources/stacks-snapshot-20261011/tags.txt, source lines 5152--5152. Navigation pointer into the Stacks Project; the standard fact it supplies is applied inside the counted proof span and the argument resting on it stays there.\par\medskip
\ifdefined\hypertarget\hypertarget{sourcepointer-2}{}\fi
\noindent\textit{Source reference.} Original lemma 06RU: focused-seeds/AGP/sources/stacks-snapshot-20261011/tags.txt, source lines 8653--8653. Navigation pointer into the Stacks Project; the standard fact it supplies is applied inside the counted proof span and the argument resting on it stays there.\par\medskip
\ifdefined\hypertarget\hypertarget{sourcepointer-3}{}\fi
\noindent\textit{Source reference.} Original lemma 02GL: focused-seeds/AGP/sources/stacks-snapshot-20261011/tags.txt, source lines 3232--3232. Navigation pointer into the Stacks Project; the standard fact it supplies is applied inside the counted proof span and the argument resting on it stays there.\par\medskip
\ifdefined\hypertarget\hypertarget{sourcepointer-4}{}\fi
\noindent\textit{Source reference.} Original lemma 06RT: focused-seeds/AGP/sources/stacks-snapshot-20261011/tags.txt, source lines 8652--8652. Navigation pointer into the Stacks Project; the standard fact it supplies is applied inside the counted proof span and the argument resting on it stays there.\par\medskip
\ifdefined\hypertarget\hypertarget{sourcepointer-5}{}\fi
\noindent\textit{Source reference.} Original lemma 01TH: focused-seeds/AGP/sources/stacks-snapshot-20261011/tags.txt, source lines 2374--2374. Navigation pointer into the Stacks Project; the standard fact it supplies is applied inside the counted proof span and the argument resting on it stays there.\par\medskip
\ifdefined\hypertarget\hypertarget{sourcepointer-6}{}\fi
\noindent\textit{Source reference.} Original lemma 03FU: focused-seeds/AGP/sources/stacks-snapshot-20261011/tags.txt, source lines 4505--4505. Navigation pointer into the Stacks Project; the standard fact it supplies is applied inside the counted proof span and the argument resting on it stays there.\par\medskip
\ifdefined\hypertarget\hypertarget{sourcepointer-7}{}\fi
\noindent\textit{Source reference.} Original lemma 03H4: focused-seeds/AGP/sources/stacks-snapshot-20261011/tags.txt, source lines 4553--4553. Navigation pointer into the Stacks Project; the standard fact it supplies is applied inside the counted proof span and the argument resting on it stays there.\par\medskip
\begin{lemma}
\label{lemma-locally-quasi-finite}
Let $S$ be a scheme. Let $f : X \to Y$ be a morphism of algebraic spaces.
Assume $f$ is locally of finite type. The following are equivalent
\begin{enumerate}
\item $f$ is locally quasi-finite,
\item for every morphism $\Spec(k) \to Y$ where $k$ is a field
the space $|X_k|$ is discrete. Here $X_k = \Spec(k) \times_Y X$.
\end{enumerate}
\end{lemma}

\begin{proof}
Assume $f$ is locally quasi-finite. Let $\Spec(k) \to Y$ be as
in (2). Choose a surjective \'etale morphism $U \to X$ where $U$ is a scheme.
Then $U_k = \Spec(k) \times_Y U \to X_k$ is an
\'etale morphism of algebraic spaces by
Properties of Spaces, Lemma \ref{spaces-properties-lemma-base-change-etale}.
By
Lemma \ref{lemma-base-change-quasi-finite}
we see that $X_k \to \Spec(k)$ is locally quasi-finite. By
definition this means that $U_k \to \Spec(k)$ is locally
quasi-finite. Hence $|U_k|$ is discrete by
Morphisms, Lemma \ref{morphisms-lemma-locally-quasi-finite-fibres}.
Since $|U_k| \to |X_k|$ is surjective and open we conclude that $|X_k|$
is discrete.

\medskip\noindent
Conversely, assume (2). Choose a surjective \'etale morphism $V \to Y$
where $V$ is a scheme. Choose a surjective \'etale morphism
$U \to V \times_Y X$ where $U$ is a scheme.
Note that $U \to V$ is locally of finite type as $f$ is locally of finite type.
Picture
$$
\xymatrix{
U \ar[r] \ar[rd] & X \times_Y V \ar[d] \ar[r] & V \ar[d] \\
& X \ar[r] & Y
}
$$
If $f$ is not locally quasi-finite then $U \to V$ is not locally quasi-finite.
Hence there exists a specialization $u \leadsto u'$ for some $u, u' \in U$
lying over the same point $v \in V$, see
Morphisms, Lemma \ref{morphisms-lemma-quasi-finite-at-point-characterize}.
We claim that $u, u'$ do not have the same image in
$X_v = \Spec(\kappa(v)) \times_Y X$
which will contradict the assumption that $|X_v|$ is discrete as desired.
Let $d = \text{trdeg}_{\kappa(v)}(\kappa(u))$ and
$d' = \text{trdeg}_{\kappa(v)}(\kappa(u'))$.
Then we see that $d > d'$ by
Morphisms, Lemma \ref{morphisms-lemma-dimension-fibre-specialization}.
Note that $U_v$ (the fibre of $U \to V$ over $v$) is the fibre product of
$U$ and $X_v$ over $X \times_Y V$, hence $U_v \to X_v$ is \'etale (as a
base change of the \'etale morphism $U \to X \times_Y V$). If
$u, u' \in U_v$ map to the same element of $|X_v|$ then there exists a point
$r \in R_v = U_v \times_{X_v} U_v$ with $t(r) = u$ and $s(r) = u'$, see
Properties of Spaces, Lemma \ref{spaces-properties-lemma-points-cartesian}.
Note that $s, t : R_v \to U_v$ are \'etale morphisms of schemes
over $\kappa(v)$, hence $\kappa(u) \subset \kappa(r) \supset \kappa(u')$
are finite separable extensions of fields over $\kappa(v)$ (see
Morphisms, Lemma \ref{morphisms-lemma-etale-over-field}).
We conclude that the transcendence degrees are equal.
This contradiction finishes the proof.
\end{proof}

\end{document}
